\documentclass[10pt,a4paper]{amsart}
\usepackage[margin=19mm]{geometry}
\usepackage{amsmath,amssymb,mathtools}
\usepackage[scr=rsfs,cal=euler]{mathalfa}
\usepackage{xfrac}
\usepackage[unicode,hidelinks,bookmarks=false]{hyperref}
\newcommand{\ie}{i.e.\ }
\newcommand{\cf}{cf.\ }
\newcommand{\resp}{respectively\ }
\newcommand{\Subsection}{\subsection}
\providecommand{\nobreakdash}{\nobreak\hbox{-}}
\setlength{\emergencystretch}{2em}
\multlinegap0pt
\usepackage{bm}
\usepackage{bbm}
\usepackage{dsfont}
\usepackage{xspace}
\usepackage{cancel}
\usepackage{mathtools}
\usepackage[shortlabels]{enumitem}
\usepackage{amsthm}
\makeatletter
\@ifundefined{theorem}{\newtheorem{theorem}{Theorem}}{}
\makeatother
\newcommand{\Aff}{\mathrm{Aff}}
\newcommand{\R}{\mathbb{R}}
\newcommand{\ex}{\mathrm{pro}}
\newcommand{\id}{\mathrm{Id}}
\title[Noncommutative protori and inductive spectral triples]{Noncommutative protori and inductive spectral triples}
\author{Remus Floricel, Patrick Melanson}
\date{Source: arXiv:2605.26049v1 (CC BY 4.0)}
\begin{document}
\maketitle

\noindent\emph{Source and licence.} Noncommutative protori and inductive spectral triples. Version arXiv:2605.26049v1 (CC BY 4.0), distributed under the Creative Commons Attribution 4.0 International licence, as recorded in the arXiv licence metadata for this exact version and shown on the abstract page https://arxiv.org/abs/2605.26049v1. Full rights notice: licenses/arxiv-2605.26049-CC-BY-4.0.txt.\par\smallskip

\noindent\textit{Editorial scope.} Counted unit: the Theorem whose source label is \texttt{thm:protorus-invariant} in the article (source0.tex lines 1515--1552, source label \texttt{thm:protorus-invariant}), delivered together with the article's own complete original proof of it (source0.tex lines 1554--1796, one \texttt{proof} environment of 243 lines). No mathematics is written, repaired or re-derived: statement and proof are the article's own text line for line. Required context retained uncounted: none. Every definition, display and label of those lines is retained unchanged. Omitted from this package and not redistributed: 1--1514, 1553--1553, 1797--6546. Nothing inside a retained excerpt is omitted or altered: every retained body is delivered exactly as the article's lines are written, so no omission receipt of any kind accompanies this package. Citations: the retained excerpts carry no citation command at all, so no bibliography entry of the article is transcribed and no reference is redistributed. Reference adaptation: none was needed and none was made; every reference inside the delivered unit resolves inside it, so no unresolved reference or citation is produced. Printed slips kept literal: no printed word, symbol, hypothesis or number of the counted statement or of its proof was altered. Layout and class: the article's own class is \texttt{amsart} with packages including geometry, amsmath, amssymb, amsfonts, amsthm, amsrefs, bbm, mathrsfs, hyperref, enumitem, mathtools, amscd, tikz-cd; the wrapper is the lane's pinned \texttt{amsart} editorial wrapper at 10pt on A4, which restates the theorem-like environments the retained text needs and adds the article's own macro definitions that the retained text needs (\texttt{\textbackslash Aff}, \texttt{\textbackslash R}, \texttt{\textbackslash ex}, \texttt{\textbackslash id}), with extra wrapper packages (bm, bbm, dsfont, xspace, cancel, mathtools, enumitem). The delivered numbering is the unit's own. Assets: no retained span contains a figure environment, a graphics-inclusion command, a PDF-page inclusion, a table or a TikZ picture; no image file is copied into this package, the package declares no asset dependency and no image file is redistributed with it. Licence: version arXiv:2605.26049v1 of this article is distributed under the Creative Commons Attribution 4.0 International licence, as recorded in the arXiv licence metadata for this exact version and shown on the abstract page https://arxiv.org/abs/2605.26049v1. A full rights notice is delivered at licenses/arxiv-2605.26049-CC-BY-4.0.txt. Characters: every retained excerpt is pure ASCII, so the delivered engine, XeLaTeX with its default Latin Modern text font, reports no missing character.
\begin{theorem}[Elliott invariant of a protoral system]
\label{thm:protorus-invariant}
Let $A_{\ex}=\varinjlim(B_n,\phi_n)$
be as above. 
\begin{enumerate}[label={\rm(\roman*)}]
\item For \(i=0,1\), one has $K_i(A_{\ex})\cong \varinjlim(K_i(B_n),f_{i,n}).$ The positive cone on \(K_0(A_{\ex})\) is the direct-limit cone.

\item The cone \(\widetilde T(A_{\ex})\) of densely defined lower semicontinuous
traces on \(A_{\ex}\) is one-dimensional. More precisely, there is a densely
defined lower semicontinuous trace \(\tau\) on \(A_{\ex}\), unique up to
multiplication by a positive scalar, such that $\tau\circ\iota_n=c_n\tau_n$
%\qquad
for all $n\geq 1$.
Equivalently, if $\rho_\tau:=\tau_*\colon K_0(A_{\ex})\to\mathbb R$
then $\rho_\tau\circ(\iota_n)_*=c_n\rho_n$ for all $n\geq 1$.
The full trace pairing $\rho_{A_{\ex}}\colon K_0(A_{\ex})\to \Aff(\widetilde T(A_{\ex}))$
is given by $\rho_{A_{\ex}}(x)(\lambda\tau)=\lambda\,\rho_\tau(x)$ for all
$x\in K_0(A_{\ex})$ and $\lambda\in[0,\infty).$

\item The scale function is $\Sigma_{A_{\ex}}(\lambda\tau)
=
\lambda\,\sup_{n\geq 1}c_n,$ for every $
\lambda\in[0,\infty),$
where the value \(+\infty\) is allowed, with the convention $0\cdot\infty=0$.

\item The projection scale is $\Sigma(K_0(A_{\ex}))
=
\bigcup_{n\geq 1}
[0,(\iota_n)_*([1_{B_n}])]
\subseteq K_0(A_{\ex})^+.$

\item If every $\phi_n$ is unital, then $A_{\ex}$ is unital,
$t_n=1$ for all $n$, and $\tau$ is the unique tracial state on
$A_{\ex}$.

\item If \(A_{\ex}\) is nonunital and \(c_n\to\infty\), then $\Sigma(K_0(A_{\ex}))=K_0(A_{\ex})^+.$
\end{enumerate}
\end{theorem}

\begin{proof}
Set $e_n:=\iota_n(1_{B_n})\in A_{\ex}.$ Then $(e_n)_{n\geq1}$ is an increasing approximate unit of projections
for $A_{\ex}$.
\smallskip

\noindent
\textup{(i)}
Continuity of $K$-theory for inductive limits gives $K_i(A_{\ex})\cong \varinjlim(K_i(B_n),f_{i,n})$ for 
$i=0,1.$ It remains to identify the positive cone. Let $V(A)$ denote the
Murray--von Neumann semigroup of projections in matrix algebras over a
$C^*$-algebra $A$. The functor $V$ is continuous for inductive limits, so $V(A_{\ex})\cong\varinjlim V(B_n).$
Each $B_n$ has real rank zero and stable rank one. Hence the natural map $V(B_n)\longrightarrow K_0(B_n)^+$
is an isomorphism of ordered semigroups. Passing to the direct limit gives that
$K_0(A_{\ex})^+$ is precisely the direct-limit cone.

\smallskip

\noindent
\textup{(ii)}
First we construct the trace. Since $\tau_{n+1}\circ\phi_n=t_n\tau_n,$ the traces $c_n\tau_n$
are compatible with the connecting maps: $c_{n+1}\tau_{n+1}\circ\phi_n
=c_{n+1}t_n\tau_n
=c_n\tau_n.$
Thus they define a positive trace $\tau_0$ on the algebraic inductive limit $\bigcup_{n\geq1}\iota_n(B_n),$ defined by
$\tau_0(\iota_n(a)):=c_n\tau_n(a),$ if
$a\in B_n.$
Let $\tau$ be its lower semicontinuous regularization on the $C^*$-completion
$A_{\ex}$:
\[
\tau(x):=\sup\{\tau_0(b):\, b\in (\bigcup_n \iota_n(B_n))_+, \ b\le x\}.
\]
Then $\tau$ is a densely defined lower semicontinuous trace on $A_{\ex}$,
and, since each $c_n\tau_n$ is bounded on $B_n$, the regularization agrees
with $c_n\tau_n$ on the finite stages, i.e.  $\tau\circ\iota_n=c_n\tau_n$ for all
$n\geq1$.
Indeed, if $x\in \iota_n(B_n)_+$ and $b\in\bigcup_m\iota_m(B_m)_+$ satisfies
$b\leq x$, then, after passing to a common later stage, the compatibility of
the traces gives $\tau_0(b)\leq \tau_0(x)$. Hence the regularization takes the
value $\tau_0(x)=c_n\tau_n(x)$ on $x$.

Passing to $K_0$ gives $\rho_\tau\circ(\iota_n)_*=c_n\rho_n,$
where $\rho_\tau=\tau_*$. This proves the displayed formula.

We now prove uniqueness of the trace ray. For $m\geq n$, write $\phi_{m,n}:=\phi_{m-1}\circ\cdots\circ\phi_n$ and
$\phi_{n,n}:=\id_{B_n}.$ Inside \(B_m\), put $p_{n,m}:=\phi_{m,n}(1_{B_n})$, and set $C_{n,m}:=p_{n,m}B_m p_{n,m}$ for simplicity.
Then $C_{n,m}$ is a nonzero full corner of the simple algebra $B_m$. Hence
$C_{n,m}$ is simple and has a unique tracial state, namely
\[
\omega_{n,m}(x)
=
\frac{\tau_m(x)}{\tau_m(p_{n,m})},
\qquad x\in C_{n,m}.
\]
The maps $C_{n,m}\longrightarrow C_{n,m+1}$
induced by $\phi_m$ are unital. Moreover, the traces $\omega_{n,m}$ are
compatible, because for $x\in C_{n,m}$,
\[
\frac{\tau_{m+1}(\phi_m(x))}
     {\tau_{m+1}(\phi_m(p_{n,m}))}
=
\frac{t_m\tau_m(x)}
     {t_m\tau_m(p_{n,m})}
=
\omega_{n,m}(x).
\]
Therefore $e_nA_{\ex}e_n
\cong
\varinjlim_{m\geq n}(C_{n,m},\phi_m|_{C_{n,m}})$
is a unital $C^*$-algebra with a unique tracial state. Denote this state by $\omega_n.$

Let $\sigma\in\widetilde T(A_{\ex})$ be a nonzero densely defined lower
semicontinuous trace. Since $\sigma$ is densely defined, it is finite on the
Pedersen ideal of $A_{\ex}$. In particular, it is finite on projections, so $\lambda_n:=\sigma(e_n)<\infty$ for every  $n\geq1$.
The restriction of $\sigma$ to the unital algebra $\iota_n(B_n)$, whose unit
is $e_n$, is therefore a bounded trace. Since $B_n$ has a unique tracial
state, there is a scalar $\lambda_n\geq0$ such that $\sigma\circ\iota_n=\lambda_n\tau_n.$
Compatibility with the connecting maps gives
\[
\lambda_n\tau_n
=
\sigma\circ\iota_n
=
\sigma\circ\iota_{n+1}\circ\phi_n
=
\lambda_{n+1}\tau_{n+1}\circ\phi_n
=
\lambda_{n+1}t_n\tau_n.
\]
Thus $\lambda_{n+1}={\lambda_n}{t_n}^{-1},$ and hence
$\lambda_n=\lambda_1c_n,$
for all $n\geq1$.

If $\lambda_1=0$, then $\lambda_n=0$ for every $n$. Since $e_nA_{\ex}e_n$
has unique tracial state, this implies $\sigma|_{e_nA_{\ex}e_n}=0$
for all $n\geq1$. For $a\in (A_{\ex})_+$, the positive elements $a^{1/2}e_na^{1/2}$
increase to $a$ and converge to
$a$ in norm. By traciality,
$\sigma(a^{1/2}e_na^{1/2})=\sigma(e_nae_n)=0.$
Lower semicontinuous traces satisfy monotone convergence, so $\sigma(a)=\sup_n\sigma(a^{1/2}e_na^{1/2})=0.$
Thus $\sigma=0$, contrary to our assumption. Hence $\lambda_1>0$.

Since $e_nA_{\ex}e_n$ has unique tracial state and $\sigma(e_n)=\lambda_n$, we have $\sigma|_{e_nA_{\ex}e_n}=\lambda_n\omega_n
=\lambda_1 c_n\omega_n.$
On the other hand, the trace $\tau$ constructed above satisfies $\tau(e_n)=c_n,$
and hence $\tau|_{e_nA_{\ex}e_n}=c_n\omega_n.$
Therefore $\sigma|_{e_nA_{\ex}e_n}=\lambda_1\,\tau|_{e_nA_{\ex}e_n}$ for every $n\geq1$.

Now let $a\in (A_{\ex})_+$. Again $a^{1/2}e_na^{1/2}\nearrow a,$
and by traciality,
\[
\sigma(a^{1/2}e_na^{1/2})
=
\sigma(e_nae_n)
=
\lambda_1\tau(e_nae_n)
=
\lambda_1\tau(a^{1/2}e_na^{1/2}).
\]
Using monotone convergence for both lower semicontinuous traces gives
\[
\sigma(a)
=
\sup_n\sigma(a^{1/2}e_na^{1/2})
=
\lambda_1\sup_n\tau(a^{1/2}e_na^{1/2})
=
\lambda_1\tau(a).
\]
Thus every nonzero densely defined lower semicontinuous trace is a positive
scalar multiple of $\tau$. Hence $\widetilde T(A_{\ex})=\R_+\tau.$ The final displayed formula in \textup{(ii)} follows from homogeneity:
for $x\in K_0(A_{\ex})$ and $\lambda\geq0$, $\rho_{A_{\ex}}(x)(\lambda\tau)=\lambda\,\tau_*(x).$

\smallskip

\noindent
\textup{(iii)}
Let $h\in (A_{\ex})_+$ be a strictly positive element. By the definition of the scale
function in the unified invariant, $\Sigma_{A_{\ex}}(\tau)=d_\tau(h).$
The support projection of $h$ in $A_{\ex}^{**}$ is $1_{A_{\ex}^{**}}$. Let
$\bar\tau$ denote the normal extension of $\tau$ to $A_{\ex}^{**}$. Then
$d_\tau(h)=\bar\tau(1_{A_{\ex}^{**}}).$
Since $(e_n)$ is an increasing approximate unit of projections, $e_n\nearrow
1_{A_{\ex}^{**}}$ strongly. Therefore
\[
\bar\tau(1_{A_{\ex}^{**}})
=
\sup_{n\geq1}\tau(e_n)
=
\sup_{n\geq1}c_n.
\]
Thus $\Sigma_{A_{\ex}}(\tau)=\sup_{n\geq1}c_n.$
By homogeneity of the scale function, $
\Sigma_{A_{\ex}}(\lambda\tau)=\lambda\,\sup_{n\geq1}c_n,$ for $
\lambda\in[0,\infty).$

\smallskip

\noindent
\textup{(iv)}
We prove the projection-scale formula $\Sigma(K_0(A_{\ex}))=\bigcup_{n\geq1}[0,[e_n]].$
Here and below $[e_n]$ means $(\iota_n)_*([1_{B_n}])$.

First let $p\in A_{\ex}$ be a projection. Since $e_n\to 1$ strictly, we have $e_npe_n\to p$ in norm. For $n$
sufficiently large, $e_npe_n$ is close enough to $p$ that the standard
projection perturbation lemma gives a projection $q\in e_nA_{\ex}e_n$, equivalently
$q\leq e_n$, which is Murray--von Neumann equivalent to $p$. Hence
$
[p]=[q]\leq [e_n].$
Therefore every projection class in $A_{\ex}$ belongs to
$
\bigcup_{n\geq1}[0,[e_n]].$

Conversely, suppose $x\in K_0(A_{\ex})^+$ and
$0\leq x\leq [e_n]$ for some $n$. Since $A_{\ex}$ is simple and $e_n\neq0$, the corner $e_nA_{\ex}e_n$
is full. Hence the inclusion $e_nA_{\ex}e_n\hookrightarrow A_{\ex}$
induces an order isomorphism on $K_0$. Let $x'\in K_0(e_nA_{\ex}e_n)^+$ be the
preimage of $x$. Then $0\leq x'\leq [1_{e_nA_{\ex}e_n}].$

The algebra $e_nA_{\ex}e_n$ has real rank zero and stable rank one, being a unital
corner of $A_{\ex}$. Therefore $x'$ is represented by a projection
in $e_nA_{\ex}e_n$.
Indeed, as $0\leq x'\leq [1_{e_nA_{\ex}e_n}]$, choose projections $q,r$ in matrix
algebras over $e_nA_{\ex}e_n$ with $[q]=x'$ and 
%\qquad
$[r]=[1_{e_nA_{\ex}e_n}]-x'.$
Then $[q]+[r]=[1_{e_nA_{\ex}e_n}].$ Since $e_nA_{\ex}e_n$ has stable rank one, projections over
$e_nA_{\ex}e_n$ satisfy cancellation; equivalently, the natural map $V(e_nA_{\ex}e_n)\to K_0(e_nA_{\ex}e_n)^+$
is injective. Hence the equality $[q]+[r]=[1_{e_nA_{\ex}e_n}]$
in $K_0$ implies that $q\oplus r$ is Murray--von Neumann equivalent to
$1_{e_nA_{\ex}e_n}$ in a matrix algebra. Therefore $q$ is Murray--von
Neumann equivalent to a subprojection of $1_{e_nA_{\ex}e_n}$, so $x'$ is
represented by an actual projection in $e_nA_{\ex}e_n$.

Therefore there exists a projection $p\in e_nA_{\ex}e_n\subseteq A_{\ex}$
such that
$[p]=x.$ Thus every element of $[0,[e_n]]$ is a projection class in $A_{\ex}$. This proves
$
\Sigma(K_0(A_{\ex}))
=\bigcup_{n\geq1}[0,[e_n]]
=\bigcup_{n\geq1}[0,(\iota_n)_*([1_{B_n}])].$

\smallskip

\noindent
\textup{(v)}
Assume every $\phi_n$ is unital. Then
$\phi_n(1_{B_n})=1_{B_{n+1}},$
so $e_n=e_{n+1}$
for all $n$. Hence $A_{\ex}$ is unital with unit $e_1$. Moreover, $t_n
=\tau_{n+1}(\phi_n(1_{B_n}))
=\tau_{n+1}(1_{B_{n+1}})=1.$
Thus $c_n=1$ for all $n$. The trace $\tau$ satisfies
$\tau(e_1)=1,$
so $\tau$ is a tracial state. Since the whole lower semicontinuous trace cone
is $\R_+\tau$, this tracial state is unique.

\smallskip

\noindent
\textup{(vi)}
Assume $A_{\ex}$ is nonunital and $c_n\to\infty$. Let
$x\in K_0(A_{\ex})^+.$
If $x=0$, then $x\in\Sigma(K_0(A_{\ex}))$.
Now assume $x\neq0$. By \textup{(i)}, there exist $n$ and $y\in K_0(B_n)^+\setminus\{0\}$ such
that $(\iota_n)_*(y)=x.$
Let $m\geq n$, and write
$f_{0,m,n}:=f_{0,m-1}\circ\cdots\circ f_{0,n}$
for the induced map from \(K_0(B_n)\) to \(K_0(B_m)\). Put
$y_m:=f_{0,m,n}(y)\in K_0(B_m)^+.$ But 
$\rho_\tau(x)=c_m\rho_m(y_m)$ by \textup{(ii)}
so $\rho_m(y_m)=c_m^{-1}\rho_\tau(x).$
Since $c_m\to\infty$, choose $m\geq n$ so large that
$c_m>\rho_\tau(x).$
Then $\rho_m(y_m)<1=\rho_m([1_{B_m}]).$
Hence $\rho_m([1_{B_m}]-y_m)>0.$ But
the order on \(K_0(B_m)\) is determined by the unique trace, so
$[1_{B_m}]-y_m\in K_0(B_m)^+.$
Thus $0\leq y_m\leq [1_{B_m}].$
Applying $(\iota_m)_*$, we get
$0\leq x\leq (\iota_m)_*([1_{B_m}])=[e_m].$
By \textup{(iv)}, $x\in\Sigma(K_0(A_{\ex})).$
Therefore $\Sigma(K_0(A_{\ex}))=K_0(A_{\ex})^+$, as claimed. The proof is complete.
\end{proof}

\end{document}
